\section{Потоки в сетях. Теорема Форда и Фалкерсона. Алгоритмы нахождения максимального потока}
\label{sec:network-flows}

\subsection{Транспортная сеть и допустимый поток}

Транспортной сетью называется ориентированный граф
\[
    G=(V,E)
\]
с выделенными вершинами $s$ --- источником и $t$ --- стоком, а также функцией
пропускных способностей
\[
    c\colon E\to\mathbb R_{\geq0}.
\]

Поток --- функция $f\colon E\to\mathbb R$, удовлетворяющая условиям:
\begin{enumerate}
    \item ограничение пропускной способности:
    \[
        0\leq f(u,v)\leq c(u,v);
    \]
    \item сохранение потока во внутренних вершинах:
    \[
        \sum_{(u,v)\in E}f(u,v)
        =
        \sum_{(v,w)\in E}f(v,w),
        \qquad v\in V\setminus\{s,t\}.
    \]
\end{enumerate}

Величина потока определяется как чистый выход из источника:
\[
 |f|=
 \sum_{(s,v)\in E}f(s,v)
 -
 \sum_{(v,s)\in E}f(v,s).
\]
По закону сохранения она равна чистому входу в сток.

\subsection{Остаточная сеть}

Для текущего потока $f$ остаточная пропускная способность прямой дуги равна
\[
    c_f(u,v)=c(u,v)-f(u,v),
\]
а обратной дуги ---
\[
    c_f(v,u)=f(u,v).
\]
Остаточная сеть $G_f$ содержит все дуги с положительной остаточной пропускной
способностью.

Путь из $s$ в $t$ в остаточной сети называется увеличивающим. Его остаточная
пропускная способность
\[
    \Delta=\min_{(u,v)\in P}c_f(u,v).
\]
Поток увеличивают на $\Delta$ вдоль прямых дуг пути и уменьшают на $\Delta$
вдоль обратных. Обратные дуги принципиальны: они позволяют отменять ранее
принятые локально неудачные решения.

\subsection{Разрез сети}

Разрезом называется разбиение вершин
\[
    V=S\cup T,
    \qquad S\cap T=\varnothing,
    \qquad s\in S,
    \qquad t\in T.
\]
Его пропускная способность равна
\[
    c(S,T)=\sum_{u\in S,\,v\in T}c(u,v).
\]
Для любого допустимого потока и любого разреза
\[
    |f|\leq c(S,T).
\]
Следовательно, пропускная способность любого разреза является верхней оценкой
максимального потока.

\subsection{Теорема Форда--Фалкерсона}

\textbf{Теорема.}
Максимальная величина потока из $s$ в $t$ равна минимальной пропускной способности
разреза:
\[
    \max_f |f|=
    \min_{(S,T)}c(S,T).
\]

\textbf{Схема доказательства.}
Если в остаточной сети нет пути из $s$ в $t$, обозначим через $S$ множество
вершин, достижимых из $s$, и положим $T=V\setminus S$. Тогда все прямые дуги из
$S$ в $T$ насыщены, а поток по всем дугам из $T$ в $S$ равен нулю. Поэтому
\[
    |f|=c(S,T).
\]
С одной стороны, поток не превосходит ёмкости любого разреза; с другой стороны,
для найденного разреза достигается равенство. Следовательно, поток максимален,
а разрез минимален.

\subsection{Алгоритм Форда--Фалкерсона}

\begin{enumerate}
    \item Инициализировать нулевой поток $f=0$.
    \item Пока в остаточной сети существует увеличивающий путь $P$:
    \begin{enumerate}
        \item найти $\Delta=\min_{e\in P}c_f(e)$;
        \item увеличить поток вдоль $P$ на $\Delta$.
    \end{enumerate}
    \item Если увеличивающего пути нет, вернуть поток $f$.
\end{enumerate}

При целых пропускных способностях каждое увеличение повышает величину потока не
менее чем на единицу, поэтому алгоритм завершается. Оценка времени при произвольном
поиске пути:
\[
    O(mF),
\]
где $F$ --- величина максимального целочисленного потока. Эта оценка
псевдополиномиальна. Для иррациональных ёмкостей неудачный выбор путей может даже
привести к незавершению.

\subsection{Алгоритм Эдмондса--Карпа}

В алгоритме Эдмондса--Карпа увеличивающий путь ищется с помощью BFS.
Поэтому выбирается путь с минимальным числом дуг. Его сложность
\[
    O(nm^2).
\]
Эта оценка полиномиальна и не зависит от величин пропускных способностей.

\subsection{Алгоритм Диница}

Алгоритм Диница строит слоистую сеть BFS и затем находит в ней блокирующий поток
с помощью DFS. После блокирующего потока расстояние от $s$ до $t$ в остаточной
сети увеличивается. Для общего графа сложность составляет
\[
    O(n^2m).
\]
На практике алгоритм Диница часто значительно быстрее Эдмондса--Карпа.

\subsection{Целочисленность и применения}

Если все пропускные способности целые, существует максимальный поток с целыми
значениями на дугах. Это позволяет сводить к потоку задачи о паросочетаниях,
назначениях, распределении ресурсов и непересекающихся путях.

Для двудольного графа максимальное паросочетание строится сетью с единичными
пропускными способностями: источник соединяется с левой долей, рёбра графа идут
слева направо, правая доля соединяется со стоком.

\subsection{Что важно сказать на экзамене}

Ключевые понятия: допустимый поток, остаточная сеть, увеличивающий путь и разрез.
Теорема максимального потока и минимального разреза одновременно даёт критерий
оптимальности и обоснование алгоритма Форда--Фалкерсона.
